\section{Key Mathematical Concepts}

% \begin{frame}{Key Mathematical Concepts}
%     \begin{itemize}
%         \item \textbf{Cosine Similarity Formula} \\
%         Measures the similarity between two vectors (e.g., query and document embeddings).
        
%         \begin{center}
%         \Large
%         $$
%         \text{CosSim}(A, B) = \frac{A \cdot B}{\|A\| \, \|B\|}
%         $$
%         \end{center}
        
%         \begin{itemize}
%             \item[$A, B$:] Vector representations (e.g., query and document embeddings)  
%             \item[$A \cdot B$:] Dot product of the vectors  
%             \item[$\|A\|, \|B\|$:] Magnitude (Euclidean norm) of vectors  
%             \item[Range:] Output ranges from $-1$ (opposite) to $+1$ (identical), with $0$ = orthogonal  
%         \end{itemize}
%     \end{itemize}
%     \vspace{0.3cm}
% \end{frame}

\begin{frame}{Key Mathematical Concepts}
    \begin{itemize}
        \item \textbf{BM25 Formula} \\
        Computes the relevance score of a document $D$ for a query $Q$.
        
        \begin{center}
        \small
        $$
        \text{score}(D, Q) = \sum_{q_i \in Q} 
        IDF(q_i) \cdot 
        \frac{f(q_i, D) \cdot (k_1 + 1)}{f(q_i, D) + k_1 \cdot \left(1 - b + b \cdot \frac{|D|}{\text{avgdl}}\right)}
        $$
        \end{center}
        
        \begin{itemize}
            \item[$f(q_i, D)$:] Frequency of term $q_i$ in document $D$  
            \item[$|D|$:] Length of document $D$  
            \item[$\text{avgdl}$:] Average document length in the corpus  
            \item[$IDF(q_i)$:] Inverse document frequency of term $q_i$  
            \item[$k_1, b$:] Hyperparameters controlling term frequency saturation and length normalization
        \end{itemize}
    \end{itemize}
    \vspace{0.3cm}
\end{frame}

\begin{frame}{Key Mathematical Concepts...}
    \begin{itemize}
        \item \textbf{IDF Formula} \\
        Measures how important a term is in the whole corpus.
        
        \begin{center}
        \Large
        $$
        IDF(q_i) = \log \left( \frac{N - n(q_i) + 0.5}{n(q_i) + 0.5} + 1 \right)
        $$
        \end{center}
        
        \begin{itemize}
            \item[$N$:] Total number of documents in the corpus  
            \item[$n(q_i)$:] Number of documents containing term $q_i$  
            \item[+1:] Added for smoothing to avoid zero or negative values
        \end{itemize}
    \end{itemize}
    \vspace{0.3cm}
\end{frame}

\begin{frame}{Key Mathematical Concepts...}
    \begin{itemize}
        \item \textbf{RRF Formula} \\
        Combines rankings from multiple retrieval methods to produce a final fused score.
        
        \begin{center}
        \Large
        $$
        \textbf{RRF}(d) = \sum_{i=1}^{n} \frac{1}{k + rank_i(d)}
        $$
        \end{center}
        
        \begin{center}
        \begin{tabular}{rl}
             \textcolor{blue}{$d$ :} & Candidate document \\
             \textcolor{blue}{$rank_i(d)$ :} & Rank of document $d$ in retrieval list $i$ \\
             \textcolor{blue}{$n$ :} & Total number of retrieval lists being fused \\
            \textcolor{blue}{$k$ :} & Constant to dampen the effect of top ranks \\
             \textcolor{blue}{Interpretation :} & Higher $\textbf{RRF}(d)$ = higher relevance after fusion
        \end{tabular}
        \end{center}
    \end{itemize}
    \vspace{0.3cm}
\end{frame}

% \begin{frame}{Key Mathematical Concepts}
%     \begin{itemize}
%         \item \textbf{Deduplication Process} \\
%         Ensures that multiple highly similar documents are not returned redundantly.
        
%         \begin{center}
%         \begin{tabular}{p{2.5cm} p{8cm}} % left column fixed, right column wraps
%             \textcolor{blue}{$d$} : & Candidate document being considered for inclusion \\
%             \textcolor{blue}{$R'$} : & Set of already selected documents \\
%             \textcolor{blue}{$\text{Sim}(d, d_j)$} : & Similarity between $d$ and each $d_j$ in $R'$ (e.g., cosine similarity) \\
%             \textcolor{blue}{$\tau$} : & Similarity threshold (e.g., 0.8) above which document is considered duplicate \\
%             \textcolor{blue}{Interpretation} : & Only unique or sufficiently distinct documents are kept in the final results
%         \end{tabular}
%         \end{center}
%     \end{itemize}
%     \vspace{0.3cm}
% \end{frame}

\begin{frame}{Key Mathematical Concepts...}
    \begin{itemize}
        \item \textbf{Model Overview} \\
        BGE-M3 [11] is a versatile text embedding model developed by BAAI, distinguished by its capabilities in Multi-Functionality, Multi-Linguality, and Multi-Granularity. It supports over 100 languages and can process inputs ranging from short sentences to long documents of up to 8192 tokens.

        \begin{center}
        \begin{tabular}{p{3cm} p{7.5cm}} % left column fixed, right column wraps
            \(\color{blue}{e_q}\) : & Dense embedding vector of the query \(q\) \\
            \(\color{blue}{e_p}\) : & Dense embedding vector of the passage \(p\) \\
            \(\color{blue}{s_{\text{dense}}}\) : & Similarity score between \(e_q\) and \(e_p\) \\
            \(\color{blue}{\tau}\) : & Similarity threshold for relevance determination \\
            \(\color{blue}{\text{Sim}(e_q, e_p)}\) : & Similarity function (e.g., cosine similarity) \\
            \(\color{blue}{\text{Interpretation}}\) : & Higher \(s_{\text{dense}}\) indicates greater relevance between \(q\) and \(p\)
        \end{tabular}
        \end{center}
    \end{itemize}
    \vspace{0.3cm}
\end{frame}

